\subsection{Средняя скорость}
%%%%%%%%%%%%%%%%%%%%%%%%%%%%%%%%%%%%%%%%%%%%%%%%%%%%
%%%%%%%%%%%%%%%%%%%%%%%%%%%%%%%%%%%%%%%%%%%%%%%%%%%%
\z{Автобус выходит из пункта A и приходит расстояние $s  = 40 \,\text{км}$ до пункта B со средней скоростью $v_1  = 40\, \frac{\text{км}}{\text{ч}}$ и останавливается там на время $t = 20\, \text{мин}$. Затем он возвращается в пункт A, проходя расстояние $s$ со средней скоростью $v_2  = 60\, \frac{\text{км}}{\text{ч}}$. Найдите среднюю $v_\text{ср}$ и среднепутевую $v_\text{сп}$ скорости за все время движения автобуса.}%
{$v_\text{ср}=0;\, v_\text{сп}=40\, \frac{\text{км}}{\text{ч}}$}

%%%%%%%%%%%%%%%%%%%%%%%%%%%%%%%%%%%%%%%%%%%%%%%%%%%%

 \z{Первую треть времени точка движется со скоростью $v_1$, вторую треть – со скоростью $v_2$, последнюю – со скоростью $v_3$. Найдите среднюю скорость точки за время движения.}%
{$v_\text{ср}=\frac{v_1+v_2+v_3}{3}$}
	
%%%%%%%%%%%%%%%%%%%%%%%%%%%%%%%%%%%%%%%%%%%%%%%%%%%%
\z{Первую треть пути точка движется со скоростью $v_1$, вторую треть – со скоростью $v_2$, последнюю – со скоростью $v_3$. Найдите среднюю скорость точки за время движения.}%
{$v_\text{ср}=\frac{3}{\frac{1}{v_1}+\frac{1}{v_2}+\frac{1}{v_3}}$}

%%%%%%%%%%%%%%%%%%%%%%%%%%%%%%%%%%%%%%%%%%%%%%%%%%%%	
\z{Тело совершает в плоскости $xOy$ два последовательных, одинаковых по длине по перемещения со скоростями $v_1  = 20\, \frac{\text{м}}{\text{с}}$ под углом $\alpha_1  = 60\degree$ к направлению оси $Ox$ и $v_2  = 40\, \frac{\text{м}}{\text{c}}$ под углом $\alpha_2  = 120\degree$ к тому же направлению. Найдите среднюю скорость движения $v_\text{ср}$.}%
{$v_\text{ср}=23\, \frac{\text{м}}{\text{с}}$}

%%%%%%%%%%%%%%%%%%%%%%%%%%%%%%%%%%%%%%%%%%%%%%%%%%%%
\z {Первую половину пути автомобиль двигался со скоростью $v_1$, а вторую половину пути следующим образом: половину времени, оставшегося на прохождение этой половины пути, он ехал со скоростью $v_2$, а конечный отрезок всего пути с такой скоростью, что она оказалась равной средней скорости движения на первых двух участках. Чему равна средняя скорость $v_\text{ср}$ автомобиля на всем пути? Автомобиль движется прямолинейно в одном направлении.}%
{$v_\text{ср}=\frac{\sqrt{v_2(v_2+8v_1)}-v_2}{2}$}
%

